\documentclass{erp-report}

\usepackage[english]{babel}
\usepackage{lipsum}
\addto\captionsenglish{\renewcommand{\bibname}{References}}

\reporttitle{Integrated Subsurface Investigation}
\reportsubtitle{Illustrative findings for a proposed field research site}
\reportnumber{ERP-EX-2026-001}
\reportversion{1.0}
\reportdate{10 August 2026}
\reporttype{Investigation report}
\reportstatus{Final example}
\reportconfidentiality{Client issue - illustrative content}

\clientname{Example Land and Water Partnership}
\clientcontact{Dr Alex Morgan, Programme Director}

\projectname{North Field Baseline Investigation}
\projectlocation{Example District, United Kingdom}
\projectnumber{ERP-P-2026-014}

\preparedby{Jordan Lee, Project Geoscientist}
\reviewedby{Samira Patel, Technical Lead}
\approvedby{Taylor Evans, Programme Director}

\begin{document}

\makecover

\reportfrontmatter
\makedocumentcontrol

\begin{executivesummary}
  \chapterintro{%
    This example demonstrates how an Earth Rover Program report can present
    a clear technical narrative while keeping project content separate from
    document styling.
  }

  \lipsum[1]

  \begin{keyfindings}
    \begin{itemize}
      \item Three laterally persistent subsurface units were identified.
      \item The central survey corridor contains the most variable ground.
      \item No immediate constraint was identified in the illustrative data.
    \end{itemize}
  \end{keyfindings}

  \begin{recommendations}[Priority actions]
    \begin{enumerate}
      \item Confirm the interpreted boundaries with targeted sampling.
      \item Retain the baseline data for comparison with later surveys.
      \item Review the investigation scope before detailed design begins.
    \end{enumerate}
  \end{recommendations}
\end{executivesummary}

\tableofcontents
\listoffigures
\listoftables

\reportmainmatter

\chapter{Introduction}
\chapterintro{%
  This chapter defines the commission, site context, reporting basis, and
  limitations that frame the interpretation presented in later chapters.
}

\section{Background}
\lipsum[2-3]

\section{Objectives}
The illustrative investigation was designed to:
\begin{itemize}
  \item establish a reproducible baseline for the project area;
  \item identify broad lateral and vertical changes in ground conditions;
  \item highlight locations where focused follow-up work may be useful; and
  \item communicate uncertainty in a form suitable for client decisions.
\end{itemize}

\subsection{Reporting basis}
\lipsum[4][1-4]

\begin{reportnote}[Use of example content]
  Names, measurements, interpretations, and citations in this document are
  placeholders. Replace them with verified project information before issue.
\end{reportnote}

\chapter{Investigation methodology}
\chapterintro{%
  The workflow combines desk study, field observations, survey acquisition,
  quality control, and integrated interpretation.
}

\section{Desk study}
\lipsum[5]

\section{Field investigation}
\lipsum[6]

\subsection{Survey layout}
\lipsum[7]

\begin{reportfigure}
  \figureplaceholder[Project location and survey layout]{72mm}
  \caption{Placeholder for a project location plan with survey lines,
    investigation points, a north arrow, scale, and coordinate reference.}
  \label{fig:survey-layout}
\end{reportfigure}

\subsection{Acquisition sequence}
The principal field stages are summarized below.
\begin{enumerate}
  \item Confirm access, services information, and task controls.
  \item Establish survey control and record site conditions.
  \item Acquire the planned measurements with daily quality checks.
  \item Archive raw data and complete the field record.
\end{enumerate}

\section{Data quality and limitations}
\lipsum[8]

\begin{reportwarning}[Interpretive limitation]
  The plotted boundaries represent an interpretation of indirect
  measurements. They should not be treated as verified material contacts
  without appropriate intrusive confirmation.
\end{reportwarning}

\chapter{Results}
\chapterintro{%
  Results are organized by the observations most relevant to the project
  objectives, with supporting detail retained in the appendices.
}

\section{Overview}
\lipsum[9]

\begin{reporttable}
  \caption{Illustrative summary of interpreted subsurface units.}
  \label{tab:units}
  \begin{tabularx}{\textwidth}{@{}p{24mm}p{32mm}X p{25mm}@{}}
    \toprule
    \textcolor{ERPTopSoil}{\MakeUppercase{Unit}} &
    \textcolor{ERPTopSoil}{\MakeUppercase{Depth range}} &
    \textcolor{ERPTopSoil}{\MakeUppercase{Interpretation}} &
    \textcolor{ERPTopSoil}{\MakeUppercase{Confidence}} \\
    \midrule
    A & 0.0--0.8 m & Variable near-surface material & Moderate \\
    B & 0.8--3.5 m & Laterally persistent intermediate unit & High \\
    C & Below 3.5 m & Dense or competent basal material & Moderate \\
    \bottomrule
  \end{tabularx}
\end{reporttable}

\section{Spatial patterns}
Figure~\ref{fig:result-section} is a placeholder for a processed section.
\lipsum[10]

\begin{reportfigure}
  \figureplaceholder[Processed profile with interpreted boundaries]{82mm}
  \caption{Illustrative processed profile. A final figure should include
    units, axes, scale, orientation, processing notes, and uncertainty.}
  \label{fig:result-section}
\end{reportfigure}

\subsection{Near-surface unit}
\lipsum[11]

\subsection{Intermediate unit}
\lipsum[12]

\subsubsection{Local variation}
\lipsum[13]

\begin{keyfindings}[Result highlights]
  \begin{itemize}
    \item Unit B is the most laterally continuous interpreted feature.
    \item The largest local change occurs between illustrative stations 18
      and 24.
    \item Confidence decreases near both ends of the survey coverage.
  \end{itemize}
\end{keyfindings}

\chapter{Interpretation and discussion}
\chapterintro{%
  The interpretation considers the evidence as a whole and distinguishes
  observations from project implications.
}

\section{Integrated interpretation}
\lipsum[14-15]

As shown in Table~\ref{tab:units}, the interpreted sequence is intentionally
simple. The project brief \cite{projectbrief} and field record
\cite{fieldrecord} would normally provide the contextual evidence used to
test this model.

\section{Uncertainty}
\lipsum[16]

\begin{summarybox}[Interpretive confidence]
  Confidence is highest where independent observations agree and coverage is
  continuous. It is lower at survey margins and where the response may have
  more than one plausible explanation.
\end{summarybox}

\section{Implications for the project}
\lipsum[17]

\chapter{Conclusions and recommendations}
\chapterintro{%
  The conclusions answer the stated objectives; the recommendations identify
  proportionate next steps for reducing uncertainty.
}

\section{Conclusions}
\begin{enumerate}
  \item The example workflow provides continuous coverage of the defined
    survey corridor.
  \item Three broad units provide a reasonable first-order explanation of
    the placeholder data.
  \item A localized change in the central corridor warrants confirmation if
    it affects the client's next decision.
  \item The stated limitations should remain attached to any reuse of the
    interpreted outputs.
\end{enumerate}

\section{Recommendations}
\begin{recommendations}
  \begin{enumerate}
    \item Replace all placeholder material with verified project evidence.
    \item Confirm the central anomaly using a targeted independent method.
    \item Review the conceptual model after new evidence becomes available.
    \item Record any departure from the stated scope in the next revision.
  \end{enumerate}
\end{recommendations}

\reportbackmatter
\addcontentsline{toc}{chapter}{References}
\begin{thebibliography}{9}
  \bibitem{projectbrief}
  Earth Rover Program.
  \emph{North Field baseline investigation: example project brief}.
  Illustrative unpublished project document, 2026.

  \bibitem{fieldrecord}
  Earth Rover Program.
  \emph{North Field baseline investigation: example field record}.
  Illustrative unpublished field document, 2026.
\end{thebibliography}

\appendix
\chapter{Supporting field record}
\chapterintro{%
  Appendices retain supporting information without interrupting the main
  decision narrative.
}

\section{Illustrative field log}
\lipsum[18]

\begin{reporttable}
  \caption{Example field activity record.}
  \label{tab:field-log}
  \begin{tabularx}{\textwidth}{@{}p{24mm}p{30mm}X@{}}
    \toprule
    \textcolor{ERPTopSoil}{\MakeUppercase{Time}} &
    \textcolor{ERPTopSoil}{\MakeUppercase{Activity}} &
    \textcolor{ERPTopSoil}{\MakeUppercase{Record}} \\
    \midrule
    08:15 & Site check & Access and controls confirmed \\
    09:05 & Calibration & Reference response within tolerance \\
    12:40 & Quality check & Repeat line accepted \\
    16:20 & Close-out & Raw files and notes archived \\
    \bottomrule
  \end{tabularx}
\end{reporttable}

\chapter{Processing record}
\section{Illustrative configuration}
The technical listing environment can retain short, reproducible settings
without introducing a separate visual style.

\begin{technicalcode}
project: ERP-P-2026-014
coordinate_reference: REPLACE_WITH_PROJECT_CRS
sampling_interval_m: 0.50
quality_control:
  repeat_threshold_percent: 5
  retain_raw_data: true
\end{technicalcode}

\lipsum[19]

\end{document}
